\section{Correspondence between results, Lean theorems and scripts}\label{sec:correspondencia}
This appendix lists every numbered result of the paper (theorems, propositions, lemmas, corollaries, and the remarks that carry a claim) with its status labels (Section~\ref{sec:status}) and the Lean theorems or scripts that check it (Tables~\ref{tab:corr-1}, \ref{tab:corr-2} and~\ref{tab:corr-3} for Sections~\ref{sec:obligacion}, \ref{sec:categorica} and~\ref{sec:debil}, and Table~\ref{tab:corr-4} for Sections~\ref{sec:gauss} and~\ref{sec:via-futura}); ``paper proof'' means that no Lean file and no script checks it. Table~\ref{tab:corr-5} does the same for the statements that are checked by Lean or by a script but are not numbered results, so that each of the $57$ theorems of the Lean files without Mathlib whose axioms are printed in Appendix~\ref{sec:archivos}, each of the seven modules of the Lean project with Mathlib, through its main theorems, and each of the eleven scripts appears in at least one of these tables. Table~\ref{tab:corr-6} lists the definitions and hypotheses of Section~\ref{sec:obligacion}, and the two premises SEL and IPS$_K$ defined in later sections, with the places where each is used.

Lean theorems are named as in the six Lean files without Mathlib and the seven modules of the Lean project with Mathlib (Table~\ref{tab:files}). The theorems of the Mathlib project are written with their namespaces: \texttt{QAudit} (Lemma~\ref{lem:J} and Theorem~\ref{thm:cheeger}), \texttt{QAuditScaling} (Proposition~\ref{prop:scaling}), \texttt{QAuditBridge} and \texttt{QAuditScaling.\allowbreak AuditInterval}, with the examples that check non-vacuity in \texttt{QAudit.\allowbreak Audit} and \texttt{QAuditScaling.\allowbreak Audit}. The developments \texttt{QAudit} and \texttt{QAuditScaling} define $\mathfrak B_\pi$, $J$, the support graph and $w_{\min}$ independently; \texttt{QAuditBridge.\allowbreak J\_eq}, \texttt{QAuditBridge.\allowbreak supportGraph\_eq} and \texttt{QAuditBridge.\allowbreak wmin\_eq} prove that the two definitions agree (for $E\succ0$ and $P\succeq0$), so that both are about the one $J$ of Lemma~\ref{lem:J}. Scripts are named without their directory (\path{python/exact}, \path{python/high_precision} or \path{python/numeric}; Table~\ref{tab:files}). In Sections~\ref{sec:obligacion}--\ref{sec:debil} the status labels are implicit (Section~\ref{sec:status}); the tables state them explicitly by the rule given there: a statement that cites a Lean name of the files without Mathlib is \stL\ in the abstract form stated next to it, its link to $J$ is Lemma~\ref{lem:J}, \stA\ \stM\ for at least two units, and the analytic statements are \stA, with \stM\ where a theorem of the Mathlib project is cited. A label \stL\ means that Lean checks the logical form of a statement on abstract types or finite graphs, and a label \stM\ that Lean with Mathlib checks it over real matrices, in the scope stated in the row; no compilation decides whether a physical or phenomenal premise holds.

\begin{table}[htbp]\centering\small
\begin{tabular}{>{\raggedright\arraybackslash}p{0.27\textwidth}>{\raggedright\arraybackslash}p{0.09\textwidth}>{\raggedright\arraybackslash}p{0.54\textwidth}}
\toprule
Result & Status & Checked by\\
\midrule
Proposition~\ref{prop:provable} (provable descriptor equalities) & \stA & paper proof\\
Theorem~\ref{thm:coverage} (relational coverage) & \stA & paper proof; it is the case of Theorem~\ref{thm:impl} in which $R$ is an equivalence relation contained in $\ker\Phi$\\
Theorem~\ref{thm:continuum} (uncountably many classes in a fibre) & \stA & paper proof\\
Proposition~\ref{prop:extension} (form of any extension) & \stA & paper proof\\
Theorem~\ref{thm:threshold} (threshold reduction) & \stA & paper proof\\
Theorem~\ref{thm:impl} (entailment criterion) & \stL & \texttt{entails\_Su\_iff} (\texttt{IntentoSu.lean}, like the next four rows)\\
Theorem~\ref{thm:holgura} (slack and saturation) & \stL & \texttt{nonconstant\_iff}, \texttt{saturated}\\
Theorem~\ref{thm:dicotomia} (co-occurrence dichotomy) & \stL & \texttt{connected\_forces\_constant}; converse: \texttt{constant\_suff}\\
Theorem~\ref{thm:residuo} (residual under OI) & \stL & \texttt{oi\_residual}\\
Theorem~\ref{thm:robusta} (robust form) & \stL & \texttt{robust\_iff\_relational}\\
Proposition~\ref{prop:seg} (SEG and $S_u$) & \stL & (a) \texttt{seg\_implies\_Su}; (b) \texttt{segd\_iff} (\texttt{FiltracionCociente.lean}), where the predicate \texttt{SEGd} is Definition~\ref{def:SEG} and \texttt{SEG} is the split of part~(a)\\
\bottomrule
\end{tabular}
\caption{Numbered results of Section~\ref{sec:obligacion}, with their status and the Lean theorems that check them.}\label{tab:corr-1}
\end{table}

\begin{table}[htbp]\centering\small
\begin{tabular}{>{\raggedright\arraybackslash}p{0.27\textwidth}>{\raggedright\arraybackslash}p{0.09\textwidth}>{\raggedright\arraybackslash}p{0.54\textwidth}}
\toprule
Result & Status & Checked by\\
\midrule
Proposition~\ref{prop:funtor} (functors into a discrete category) & \stA & paper proof\\
Lemma~\ref{lem:J} (Gaussian integration and connectivity) & \stA, \stM & paper proof; for $n\ge2$, \texttt{QAudit.\allowbreak lemmaJ\_stationary}, with $P\succ0$ and $P=APA^\top+E$ given as data, and \texttt{QAudit.\allowbreak lemmaJ} and \texttt{QAudit.\allowbreak J\_eq\_zero\_iff} for every $P\succ0$ (together they give $J\ge0$); with $P$ constructed as the series $\sum_jA^jE(A^j)^\top$ when $\operatorname{spr}(A)<1$, \texttt{QAuditBridge.\allowbreak stationary\_of\_spr} and \texttt{QAuditBridge.\allowbreak lemmaJ\_spr}; the inner minimum is an infimum, attained by \texttt{QAudit.\allowbreak exists\_minimizer} and \texttt{QAudit.\allowbreak exists\_J\_minimizer}; $\det(I+Y)\ge1+\operatorname{tr}Y$ for $Y\succeq0$ is \texttt{QAudit.\allowbreak one\_add\_trace\_le\_det\_one\_add}; the convention $J=+\infty$ for one unit is not formalized; non-vacuity, $J=\tfrac12\log\tfrac54$ on a two-unit system, \texttt{QAudit.\allowbreak Audit.\allowbreak J\_ex\_eq} \stM; numerical illustration on the four networks, \path{J_conectividad.py} \stN; exact instance on $F_0$, $J=\tfrac12\log(1+k)$, \path{v12_familia_escala.py} \stS\\
Remark~\ref{rem:small} (systems with fewer than two units) & \stA & paper proof (conventions, and the invariance of $J$ under isomorphism); not formalized: in Lean an infimum over the empty set of bipartitions is $0$, so the Mathlib statements of Lemma~\ref{lem:J}, of Theorem~\ref{thm:cheeger} and of the strict parts of Proposition~\ref{prop:scaling} assume $n\ge2$\\
Proposition~\ref{prop:monoidal} (monoidal behavior) & \stA, \stL & \texttt{tensor\_no\_universal} (no hub on $X\otimes Y$); the rest is a paper proof; instances: $F_0\otimes F_0$, \path{gauss_conexiones.py} \stS, and a two-unit system tensored with a three-unit path, \path{pistas_numerico.py} \stN\\
Theorem~\ref{thm:translation} (translation criterion) & \stL & \texttt{nat\_iff\_functional} (\texttt{UnificacionCat.lean}), with $D_j$ nonempty (\texttt{Inhabited})\\
Theorem~\ref{thm:reconstruccion} (reconstruction) & \stL, \stA & criterion: \texttt{reconstruction\_iff\_injective}; (i): \texttt{blind\_spot}, \texttt{blind\_spot\_multiset}; paper: (ii), and the invariance of the degree multiset under isomorphism\\
Theorem~\ref{thm:bridge} (bridge consistency) & \stL, \stA & \texttt{bridge\_iff\_agree}; the uniqueness, which is immediate, is not part of the Lean statement\\
\bottomrule
\end{tabular}
\caption{Numbered results of Section~\ref{sec:categorica}, with their status and the Lean theorems or scripts that check them.}\label{tab:corr-2}
\end{table}

\begin{table}[htbp]\centering\small
\begin{tabular}{>{\raggedright\arraybackslash}p{0.27\textwidth}>{\raggedright\arraybackslash}p{0.09\textwidth}>{\raggedright\arraybackslash}p{0.54\textwidth}}
\toprule
Result & Status & Checked by\\
\midrule
Proposition~\ref{prop:decomposition} (extensive decomposition plus defect) & \stA & paper proof; exact instances on $F_0$, $F_{\rm mem}$ and $F_0\otimes F_0$, \path{gauss_conexiones.py} \stS; values on the networks of Table~\ref{tab:networks}, \path{pistas_numerico.py} \stN\\
Theorem~\ref{thm:cheeger} (min-cut sandwich) & \stA, \stM & paper proof; for $n\ge2$, exactly as stated, with the extreme eigenvalues and with $P\succ0$, $P=APA^\top+E$ given as data, \texttt{QAudit.\allowbreak sandwich}; with Loewner bounds $e_1I\preceq E\preceq e_2I$ and $P\preceq p_2I$ in place of the eigenvalues, \texttt{QAudit.\allowbreak sandwich\_lower} and \texttt{QAudit.\allowbreak sandwich\_upper}; with $P$ constructed as a series when $\operatorname{spr}(A)<1$, \texttt{QAuditBridge.\allowbreak sandwich\_spr}, and along the scaling $tA$, \texttt{QAuditBridge.\allowbreak sandwich\_along\_scaling}; on a two-unit system, where the lower bound is attained, \texttt{QAudit.\allowbreak Audit.\allowbreak sandwich\_ex} \stM; numerical check on the four networks and on twelve random systems (eleven non-degenerate), \path{pistas_numerico.py} \stN; exact check on $F_0$, where, for $c>0$, the lower bound is attained if and only if $s_2\le s_1$, \path{gauss_conexiones.py} \stS\\
Theorem~\ref{thm:laxa} (lax 2-cell) & \stL, \stA, \stM, \stS, \stN & (a) \texttt{universal\_connected}, with Lemma~\ref{lem:J} \stA\ \stM\ for $n\ge2$ and Remark~\ref{rem:small} \stA\ for $n=1$; (b) \texttt{n2\_equiv} (\texttt{GaussConexion.lean}), \texttt{n3\_equiv}; (c) \texttt{path\_reach\_from\_zero}, \texttt{path\_no\_universal}, \texttt{n4\_count}; the counts $0$, $0$, $15$ and $472$ for $n=2,\dots,5$ by exhaustive enumeration, \path{pistas_numerico.py} \stS; (d) paper proof \stA, with a numerical illustration on the star, \path{pistas_numerico.py} \stN\\
Theorem~\ref{thm:section} (factor through a section) & \stL & \texttt{section\_factor\_iff} (\texttt{PistasUnificacion.lean})\\
Theorem~\ref{thm:unif} (structural unification) & \stA, \stL, \stM, \stN & (i)--(iii): as the results they collect (Proposition~\ref{prop:decomposition}, Lemma~\ref{lem:J}, Proposition~\ref{prop:monoidal}, Theorems~\ref{thm:cheeger} and~\ref{thm:laxa}), with their combinatorial parts \stL, their parts checked with Mathlib \stM\ and numerical checks \stN; (iv): \texttt{square\_fails}, \texttt{values}, \texttt{no\_translation}, \texttt{blind\_spot}, \texttt{section\_factor\_iff}, \texttt{restriction\_vacuous}, \texttt{witness\_breaks\_purity} \stL, and Proposition~\ref{prop:extension} \stA\ for SEL against Rob\\
Theorem~\ref{thm:locus} (exact locus of non-commutation) & \stL, \stA, \stM & \texttt{noncomm\_locus}, \texttt{commute\_off\_band} (\texttt{FiltracionCociente.lean}), on Boolean predicates; their hypothesis, that a hub implies connectivity, is Theorem~\ref{thm:laxa}(a); that $J>0$ is connectivity on nonempty systems is Lemma~\ref{lem:J}, \stA\ \stM\ for at least two units, and Remark~\ref{rem:small} \stA\ for one unit\\
Theorem~\ref{thm:collapse} (collapse and zero-product monoid) & \stL, \stA, \stM & abstract form: \texttt{factors\_through\_collapse}, \texttt{zmul\_assoc} (with \texttt{zmul\_unit\_left}, \texttt{zmul\_unit\_right}) and \texttt{strong\_monoidal\_zero}, under its hypotheses \texttt{hunitL}, \texttt{hunitR}, \texttt{hne}, \texttt{hF0}, \texttt{hFne}, \texttt{hvan}; application to $J$ and to broadcast: Lemma~\ref{lem:J} \stA\ \stM, Theorem~\ref{thm:laxa}(a) and Proposition~\ref{prop:monoidal} \stA; instance on $F_0\otimes F_0$, \path{gauss_conexiones.py} \stS\\
Theorem~\ref{thm:selection} (variational selection in the fibres) & \stL & (a) \texttt{section\_factor\_iff}; (b) \texttt{restriction\_vacuous}; (c) \texttt{segd\_iff}, \texttt{seg\_implies\_Su}, \texttt{seg\_is\_Su}\\
Proposition~\ref{prop:conditional} (conditional unification) & \stA & paper proof, not formalized; its steps use \texttt{seg\_implies\_Su}, \texttt{section\_factor\_iff} and Theorem~\ref{thm:locus}\\
\bottomrule
\end{tabular}
\caption{Numbered results of Section~\ref{sec:debil}, with their status and the Lean theorems or scripts that check them.}\label{tab:corr-3}
\end{table}

\begin{table}[htbp]\centering\small
\begin{tabular}{>{\raggedright\arraybackslash}p{0.27\textwidth}>{\raggedright\arraybackslash}p{0.09\textwidth}>{\raggedright\arraybackslash}p{0.54\textwidth}}
\toprule
Result & Status & Checked by\\
\midrule
Proposition~\ref{prop:windows} (finite-window non-identifiability) & \stA, \stS & paper proof, general in $T$; the factorization of $\Sigma_T$, $\det\mathcal L_T=1$ and the value of $\det\Sigma_T$ for $T=0,\dots,5$, \path{v4_historias_ac.py} \stS\\
Proposition~\ref{prop:gauss-chains} (conditional constancy) & \stA, \stS, \stL & paper proof; the identities of (G7) and four exact chains, \path{v5_rkp_cadenas.py} \stS; an instance of \texttt{connected\_forces\_constant} \stL\\
Proposition~\ref{prop:gauss-plateau} (Plateau) & \stA, \stL & paper proof; only the consequence on the product domain, in the direction ``$S$ implies constancy'', is \texttt{plateau\_product} (\texttt{GaussConexion.lean}); the full support of the window laws is Proposition~\ref{prop:windows}\\
Theorem~\ref{thm:frontier} (exact frontier) & \stL & (a) \texttt{im\_iff\_edges}; (b) \texttt{constant\_on\_scc}, with \texttt{monotone\_on\_reach}; (c) \texttt{nonconstant\_if\_not\_scc} (\texttt{ViaFutura.lean}); in Lean the degrees take integer values, and the transcription to real values is not formalized (Section~\ref{sec:via-futura})\\
Proposition~\ref{prop:scaling} (monotonicity under scaling) & \stA, \stM, \stS, \stN & paper proof; (a) \texttt{QAuditScaling.\allowbreak prop\_scaling\_a}, on $\{t\ge0:t\operatorname{spr}(A)<1\}$ and without assuming $\operatorname{spr}(A)<1$; (b) \texttt{QAuditScaling.\allowbreak prop\_scaling\_b}, on $\mathcal C$ and without assuming that $W$ has zero diagonal; that $\mathcal C$ is an interval containing $0$, \texttt{QAuditScaling.\allowbreak AuditInterval.\allowbreak C\_interval\_containing\_zero} (with \texttt{QAuditScaling.\allowbreak AuditInterval.\allowbreak C\_ordConnected}); the support and the minimum cut, \texttt{QAuditScaling.\allowbreak supportGraph\_smul}, \texttt{QAuditScaling.\allowbreak supportGraph\_Acoup}, \texttt{QAuditScaling.\allowbreak wmin\_smul} and \texttt{QAuditScaling.\allowbreak wmin\_Acoup}; the covariance is the series $\sum_jM^jE(M^j)^\top$, convergent by \texttt{QAuditScaling.\allowbreak serSummable\_of\_spr\_lt\_one}, with $\operatorname{spr}$ Mathlib's spectral radius of the complexified matrix, and the inner minimum is attained, \texttt{QAuditScaling.\allowbreak exists\_minimizer}; in~(b) the Lean proof uses $Q(c_2)-Q(c_1)\succeq(c_2-c_1)^2E$, \texttt{QAuditScaling.\allowbreak Qb\_mono}, and states neither $Q'=AQ'A^\top+S$ nor identity~\eqref{eq:S-identity}; non-vacuity, \texttt{QAuditScaling.\allowbreak Audit.\allowbreak partA\_nonvacuous} and \texttt{QAuditScaling.\allowbreak Audit.\allowbreak partB\_nonvacuous} \stM; identity~\eqref{eq:S-identity}, with generic symbolic $3\times3$ matrices \stS\ (for general $n$, the computation in the proof \stA), and $Q'=AQ'A^\top+S$ on a rational instance, \path{escalado_diagonal_escalar.py} \stS; grids for (a), \path{monotonia_escalado.py}, and for (b), \path{escalado_diagonal_escalar.py} and \path{banda_monotonia.py} \stN\\
Remark~\ref{rem:fixed-diagonal} (the diagonal must scale, or be scalar) & \stN & \path{contraejemplo_diagonal_fija.py}: exact elimination of one variable and real roots to $50$ digits; not a certificate\\
Corollary~\ref{cor:band-scaling} (the band under scaling) & \stL, \stA, \stM, \stN & graph part: \texttt{path\_reach\_from\_zero}, \texttt{path\_no\_universal} \stL; the unchanged support, the strict increase of $J$ and the quadratic growth of the minimum cut, by Proposition~\ref{prop:scaling} \stA\ \stM\ (\texttt{QAuditScaling.\allowbreak supportGraph\_smul}, \texttt{QAuditScaling.\allowbreak supportGraph\_Acoup}, \texttt{QAuditScaling.\allowbreak prop\_scaling\_a}, \texttt{QAuditScaling.\allowbreak prop\_scaling\_b}, \texttt{QAuditScaling.\allowbreak wmin\_smul}, \texttt{QAuditScaling.\allowbreak wmin\_Acoup}); the corollary as one statement, which combines results of the files without Mathlib and of the Mathlib project, and the case of $C_n$ are \stA; grids, and $w_{\min}(tA)=t^2w_{\min}(A)$ to machine precision, \path{monotonia_escalado.py} \stN\\
\bottomrule
\end{tabular}
\caption{Numbered results of Sections~\ref{sec:gauss} and~\ref{sec:via-futura}, with their status and the Lean theorems or scripts that check them.}\label{tab:corr-4}
\end{table}

\begin{table}[htbp]\centering\small
\begin{tabular}{>{\raggedright\arraybackslash}p{0.36\textwidth}>{\raggedright\arraybackslash}p{0.09\textwidth}>{\raggedright\arraybackslash}p{0.45\textwidth}}
\toprule
Statement & Status & Checked by\\
\midrule
Countermodels CM1--CM4 (Section~\ref{sec:intento}) & \stL & \texttt{cm1\_not\_Su} (with \texttt{cm1\_OI}, \texttt{cm1\_Adm}), \texttt{cm2\_not\_Su} (with \texttt{cm2\_factor}), \texttt{cm3\_Su}, \texttt{cm3\_nonconstant}, \texttt{cm4\_not\_rel}\\
The bare blanket predicate, constant on $\mathbf{Open}$, differs from broadcast on the ring (Section~\ref{sec:proyecciones}) & \stL & \texttt{blanket\_vs\_broadcast}\\
No translation from positive integration to broadcast; none in the reverse direction (Section~\ref{sec:traduccion}) & \stL; \stA & \texttt{no\_translation}; the reverse direction follows from \texttt{values}, \texttt{nat\_iff\_functional} and Lemma~\ref{lem:J}, and is not a separate compiled statement\\
The cone of projections factors uniquely through the product (Section~\ref{sec:limite}) & \stL & \texttt{pairing\_unique}\\
The bridge square commutes on $\{K_{1,3},2K_2\}$ and not on $C_4$ or $P_4$; the values of $J$ in Table~\ref{tab:C-1} (Section~\ref{sec:puentes}) & \stL; \stN & \texttt{square\_commutes\_on\_sub}, \texttt{square\_fails}, \texttt{values}; \path{J_conectividad.py}, \path{pistas_numerico.py}\\
With OI$_\gamma$, fibre purity gives $S_u$; one non-isomorphic pair in a fibre destroys purity (Section~\ref{sec:via3}) & \stL & \texttt{purity\_gives\_Su}, \texttt{witness\_breaks\_purity}\\
Presence 2-cells into $\mathcal C_{J>0}$ are trivial (Section~\ref{sec:via4}) & \stL & \texttt{connected\_domain\_trivial}\\
On nonempty systems, $w_{\min}>0$ if and only if $J>0$ (Section~\ref{sec:l1}) & \stA, \stM & $n\ge2$: \texttt{QAuditBridge.\allowbreak J\_pos\_iff\_wmin\_pos}; $n=1$: Remark~\ref{rem:small} \stA\\
On the canonical representatives every function of the representative is a function of $\Phi$ (Section~\ref{sec:l4}) & \stL & \texttt{pullback\_determined}\\
Items (G1)--(G2) of Section~\ref{sec:gauss} & \stS & \path{v12_familia_escala.py}\\
Items (G3)--(G5) & \stS, \stA & \path{v3_modelos_apuntados.py}; written proof of (G3) \stA; for (G5) see connection (v)\\
Item (G6) & \stA, \stS & Proposition~\ref{prop:windows}; \path{v4_historias_ac.py}\\
Item (G7) & \stS & \path{v5_rkp_cadenas.py}\\
Connections (i)--(iv) of Section~\ref{sec:gauss} & \stS, \stL & \path{gauss_conexiones.py} (the analytic parts are Theorem~\ref{thm:cheeger} and Proposition~\ref{prop:gauss-plateau}); (i) \texttt{n2\_equiv}, \texttt{tensor\_no\_universal}, \texttt{f0\_hub\_readings}; (ii) \texttt{n2\_equiv}, \texttt{n3\_equiv}, \texttt{path\_reach\_from\_zero}, \texttt{path\_no\_universal}; (iii) \texttt{zmul\_assoc}, \texttt{strong\_monoidal\_zero}; (iv) \texttt{restriction\_vacuous}, \texttt{section\_factor\_iff}\\
Connection (v): the constancy results as instances; IPS$_K$ and SEG as reformulations & \stL & \texttt{connected\_forces\_constant}, \texttt{aligned\_premise\_iff}, \texttt{segd\_iff}\\
The four-unit networks: Table~\ref{tab:networks}, the hub bounds on the star and the tensor product (Section~\ref{sec:networks}) & \stN & \path{pistas_numerico.py}, \path{J_conectividad.py}\\
Exact covariance gives $S_u$ along chains and on components; under the order clause on all pairs, two classes that each occur strictly below the other have equal degrees; the $H^1$ obstruction is ill-posed (Section~\ref{sec:collapse}) & \stL; \stA & \texttt{exact\_covariance\_gives\_Su}, \texttt{exact\_covariance\_components}, \texttt{order\_all\_pairs\_general}, \texttt{order\_all\_pairs\_collapse} (integer-valued degrees in Lean); the paragraph on $H^1$ is a paper argument \stA\\
A finite model escapes the collapse; equality required on $\Phi_Q$-preserving steps that chain a fibre restores $S_u$ on it (Section~\ref{sec:im}) & \stL & \texttt{intervention\_monotonicity\_escapes}, \texttt{equality\_on\_fibre\_steps\_collapses}\\
\bottomrule
\end{tabular}
\caption{Statements checked by Lean or by a script that are not numbered results, with their location and status.}\label{tab:corr-5}
\end{table}

\begin{table}[htbp]\centering\small
\begin{tabular}{>{\raggedright\arraybackslash}p{0.30\textwidth}>{\raggedright\arraybackslash}p{0.62\textwidth}}
\toprule
Definition or hypothesis & Where it is used\\
\midrule
Definition~\ref{def:pcm} (realizations, pointed causal models, isomorphism) & Definitions~\ref{def:adm}, \ref{hyp:OI} and~\ref{hyp:Rob}; Proposition~\ref{prop:provable}; Proposition~\ref{prop:extension}, through part~(e); connection (iv) of Section~\ref{sec:gauss}\\
Definition~\ref{def:adm} (descriptors, Adm, dependence base) & Proposition~\ref{prop:provable}; the remarks after Theorem~\ref{thm:coverage} and after Definition~\ref{hyp:LSS}; Sections~\ref{sec:laguna} and~\ref{sec:intento}; Proposition~\ref{prop:funtor} (the remark after it); Proposition~\ref{prop:windows} (law-level descriptors) and connection (iv) of Section~\ref{sec:gauss}\\
Definition~\ref{def:phiQ} (Q, $D_0$, $\Phi_Q$, components, $S_x$) & every statement about $S_u$ \eqref{eq:Su}; part~(e), factorization, in Section~\ref{sec:laguna} and in Proposition~\ref{prop:conditional} (conditional unification); the architecture of Q in Section~\ref{sec:puentes}; $S_J$ in Definition~\ref{def:gauss-F0}(d); the other obligations of Q in Section~\ref{sec:mejor}\\
Definition~\ref{def:unity} (episodes, coexperience, $\mathfrak U$, $u$, $\psi$, $B_\psi$) & Definitions~\ref{hyp:UT}, \ref{hyp:pres}, \ref{hyp:N} and~\ref{def:SEG}; Theorem~\ref{thm:robusta}; Section~\ref{sec:intento}; Table~\ref{tab:B-1}; Proposition~\ref{prop:seg}; Sections~\ref{sec:l3} and~\ref{sec:l4}; connection (iv) of Section~\ref{sec:gauss}\\
Definition~\ref{def:orderB} (order clause B, A9--U, A14--U) & Theorem~\ref{thm:threshold} and the paragraph after it; Section~\ref{sec:mejor}; the discussion of Proposition~\ref{prop:gauss-chains}; Sections~\ref{sec:collapse}, \ref{sec:im} and~\ref{sec:scaling}\\
Definition~\ref{hyp:LSS} (LSS$_\Delta$) & Theorem~\ref{thm:coverage} (the remark after it); Proposition~\ref{prop:extension}(ii); Sections~\ref{sec:laguna}, \ref{sec:intento} and~\ref{sec:mejor}; the discussion of Proposition~\ref{prop:gauss-plateau}\\
Definition~\ref{hyp:OI} (OI$_\gamma$) & Theorem~\ref{thm:coverage} (the remark after it); Proposition~\ref{prop:extension}(ii); Section~\ref{sec:laguna}; the application of Theorem~\ref{thm:residuo} in Section~\ref{sec:intento}; Table~\ref{tab:B-1}; Sections~\ref{sec:mejor} and~\ref{sec:via3}; Table~\ref{tab:F-3}\\
Definition~\ref{hyp:Rob} (Rob) & Proposition~\ref{prop:extension}; Section~\ref{sec:laguna}; Table~\ref{tab:B-1}; Section~\ref{sec:mejor}; Definition~\ref{def:SEL} and Section~\ref{sec:via3}; Theorem~\ref{thm:unif}(iv); Section~\ref{sec:l3}\\
Definition~\ref{hyp:UT} (UT) & the paragraph after Theorem~\ref{thm:threshold}; Section~\ref{sec:intento}; Table~\ref{tab:B-1}; Section~\ref{sec:mejor}\\
Definition~\ref{hyp:pres} (B$_\psi^{\rm pres}$) & the remark after Definition~\ref{hyp:OI}; the paragraph after Theorem~\ref{thm:threshold}; Section~\ref{sec:intento}; Table~\ref{tab:B-1}; Section~\ref{sec:mejor}\\
Definition~\ref{hyp:N} (N) & the paragraph after Theorem~\ref{thm:threshold}; Section~\ref{sec:intento}; Table~\ref{tab:B-1}; Section~\ref{sec:mejor}\\
Definition~\ref{hyp:anchor} (absence anchor) & Theorem~\ref{thm:threshold}(v) and the paragraph after it; Section~\ref{sec:mejor}\\
Definition~\ref{def:TL} (threshold law TL) & Theorem~\ref{thm:threshold}(iii) and~(v) and the paragraph after it; Section~\ref{sec:mejor}\\
Definition~\ref{def:SEG} (SEG) & Proposition~\ref{prop:seg}; Figure~\ref{fig:graph}; Sections~\ref{sec:l3} and~\ref{sec:l4}; connection (v) of Section~\ref{sec:gauss}; Sections~\ref{sec:collapse} and~\ref{sec:im}\\
\midrule
Definition~\ref{def:SEL} (SEL, Section~\ref{sec:via3}) & Section~\ref{sec:via3}; Theorem~\ref{thm:unif}(iv); Section~\ref{sec:l3}\\
Definition~\ref{def:IPS} (IPS$_K$, Section~\ref{sec:gauss}) & item (G5); Proposition~\ref{prop:gauss-chains}; connection (v) of Section~\ref{sec:gauss}\\
\bottomrule
\end{tabular}
\caption{Definitions and hypotheses of Section~\ref{sec:obligacion}, and the two further premises defined in later sections, with the places where each is used. None of the hypotheses is adopted in this paper.}\label{tab:corr-6}
\end{table}

\clearpage
\subsection{Five candidate consequences of the categorical thesis}\label{sec:consecuencias}
One might draw five consequences from the categorical thesis of Section~\ref{sec:categorica}. Table~\ref{tab:F-3} states each one and what the results of this paper support; each entry rests only on the results and sections that it cites.
\begin{table}[htbp]\centering\small
\begin{tabular}{>{\raggedright\arraybackslash}p{0.22\textwidth}>{\raggedright\arraybackslash}p{0.70\textwidth}}
\toprule
Candidate consequence & What the results of the paper support\\
\midrule
1. The disputes between the theories come to an end & It does not follow. For the proxies studied, the framework locates the disagreement between the presence bridges of IIT type and GNWT type in squares that do not commute, exactly on the band of connected systems without a hub, such as $C_4$ and $P_4$ (Section~\ref{sec:puentes} and Theorem~\ref{thm:locus}). Whether the theories themselves disagree there is an empirical question (Section~\ref{sec:tesis}).\\
2. A dictionary between the theories given by natural transformations & Only where there is functional dependence: a translation exists if and only if the target projection is constant on the fibres of the source projection (\texttt{nat\_iff\_functional}, Section~\ref{sec:traduccion}), and there is none between positive integration and broadcast, in either direction. In general the translation is the relation $R_{ij}=\mathrm{im}\langle F_i,F_j\rangle$. The lax 2-cell of Theorem~\ref{thm:laxa} is an inequality between the presence predicates, strict for every $n\ge4$, together with lower bounds on $J$: it gives bounds, not a dictionary.\\
3. Criteria for consciousness in artificial systems & The projections give isomorphism invariants of a model (Section~\ref{sec:tesis}); which invariants are necessary for consciousness is a question about bridge laws that this paper does not settle, and Q does not fix whether $D_0$ contains artificial systems (Definition~\ref{def:phiQ}(a)). IIT and GNWT can also disagree when the grain or the substrate changes, for example on a workspace architecture implemented on a conventional computer (Section~\ref{sec:via2}, \emph{Scope}). Under OI$_\gamma$, which is not adopted, a realization isomorphic at grain $\gamma$ to a conscious one has the same total experience (Definition~\ref{hyp:OI}); so a criterion that denies experience to every artificial system is incompatible with OI$_\gamma$ as soon as some artificial system is isomorphic at grain $\gamma$ to a conscious realization. Choosing between such a criterion and OI$_\gamma$ is a decision, not a theorem.\\
4. A single clinical measure & It does not follow. The joint projection takes values in the product of the codomains of the projections, not on a single scale, and it has blind spots (Section~\ref{sec:limite} and Theorem~\ref{thm:reconstruccion}). A common scalar would need bridge laws, which this paper does not provide (strong thesis, Section~\ref{sec:tesis}), and a measurement of $u$ would need a calibrated channel, which is not available (Definition~\ref{def:orderB}(c) and Section~\ref{sec:mejor}).\\
5. Falsifiability & In principle, but not with what is available. Refuting SEG, the testable form of non-degenerate $S_u$, would need two original realizations with the same computed value of $\Phi_Q$, coexperience structures determined by a calibrated channel and a specified $\psi$, and none of the three is available (Section~\ref{sec:mejor}). The band is where the two kinds of proxy make different predictions (Section~\ref{sec:guide}). The prediction of Section~\ref{sec:scaling}, that $u$ does not decrease along the deformations of Proposition~\ref{prop:scaling}, is conditional on $\mathrm{IM}_{\mathcal T}$ with $J$ in place of $\Phi_Q$, which is not justified, and testing it needs an ordinal intra-subject measure of $u$, a calibrated measurement channel for $u$, which is not available (A14--U is a stipulated rule, not a calibration).\\
\bottomrule
\end{tabular}
\caption{Five consequences that one might draw from the categorical thesis, and what the results of this paper support.}\label{tab:F-3}
\end{table}

\paragraph{Three tempting formalizations that are vacuous.} Three ways of formalizing these consequences come to mind and prove nothing.
\begin{itemize}
\item Defining the GNWT-type predicate as $J>\theta$, for a threshold $\theta\ge0$, makes the 2-cell into the IIT-type predicate $J>0$ true by definition, since $J>\theta$ implies $J>0$. It is a higher threshold on the same integration and no longer involves broadcast, whereas the 2-cell of Theorem~\ref{thm:laxa} relates a broadcast hub to $J$.
\item A canonical section of $\Phi$ obtained from a choice function that does not use the free energy is an arbitrary choice of representatives, and no section, whatever its rule of choice, gives $S_u$: by \texttt{section\_factor\_iff} (Section~\ref{sec:via3}), the values at the representatives define a factor of $u$ if and only if $S_u$ already holds.
\item A uniqueness statement of the form ``there is exactly one $x$ with $x=c$'' holds for every $c$; it says nothing about the uniqueness of a minimizer of the free energy in a fibre, which is not guaranteed (Section~\ref{sec:via3}).
\end{itemize}
None of the three removes the uncountably many isomorphism classes that the fibre of $\Phi_Q$ through $s_0$ contains under Rob when $\theta\mapsto\Phi_Q(s_\theta)$ is continuous on the convex set $P$ of Rob and $\Phi_Q(s_0)$ is interior to its range (Theorem~\ref{thm:continuum} and Proposition~\ref{prop:extension}(i)). The formulation with content is the physical selection law SEL (Definition~\ref{def:SEL}), which contradicts Rob under the same conditions (Section~\ref{sec:via3}). Instead of these three, the paper uses the lax 2-cell of Theorem~\ref{thm:laxa}, whose combinatorial parts are compiled in Lean (Table~\ref{tab:corr-3}), and three theorems of \texttt{PistasUnificacion.lean}, \texttt{section\_factor\_iff}, \texttt{purity\_gives\_Su} and \texttt{witness\_breaks\_purity}, which state what a selection in the fibres gives and what it does not give (Section~\ref{sec:via3}).
\clearpage
